\begin{abstract}
Let $f:Y\to B$ be a torus fibration from a smooth compact complex
threefold with $b_1(Y)=b_2(Y)=0$. We identify the cotangent torsion
and connecting homomorphism omitted in the singular-fibre reduction
of the Campana--Demailly--Peternell argument, and give the exact
geometric criterion for its degree-two tangent-cohomology vanishing.
The normalization target alone does not determine that cohomology:
the explicitly computed integral twists include a one-dimensional
source whose natural map to the one-dimensional target is zero.
For irreducible and normal-crossing supports, the local conductor,
monodromy and vanishing-cycle calculations give nonnegative fibre
Euler numbers and $c_3(Y)\in\{0,2\}$. The zero-Euler irreducible
singular supports have smooth minimal elliptic ruled or Hopf
normalization. We also prove the persistence of a rank-two torus
fibration under small deformations and retain the full period,
tangent-cohomology and deformation-chain calculations. The exact
comb kernel splitting remains an explicitly stated open problem.
\end{abstract}

\paragraph{Relation to the cited results.}
The family contradicting the conclusion of
\citet[Corollary~2.3, p.~2 of the attached-source edition]{CDP20} is due to Claude and L.~Alp\"oge;
see \citet[Theorem~9.1]{AC2026} and Paper~I.
The rational component-count correction is
\cite[Section~10.5, p.~91]{AC2026}; the integral refinement below
retains its additional torsion information. The meromorphic
nonholomorphic-map theorem of
\citet[Theorem~3.1, p.~9 of the attached-source edition]{LRS2019} is used with its stated hypotheses.
The local flatness input is
\cite[Theorem~2.2, p.~5 of the attached-source edition]{Viaclovsky2026Combs}.
These page numbers refer to the attached-source editions identified
in the bibliography; their mathematical text is unchanged.
Known-result proofs remain in the S6\_09 compendium. The present
paper includes a proof only for a new result or for the precise
broader form required here.

\paragraph{The actual maps.}
A fibre denotes the full Cartier fibre, with its multiplicities.
Reduction, finite normalization, and resolution are separately
displayed morphisms. Every cohomological restriction and every
conductor comparison below refers to those morphisms. Paper~VI
proves the sign theorem for arbitrary reducible supports and
arbitrary support singularities.

\paragraph{Reading the verification record.}
Each result retains the named historical checks and separately
records a re-derivation made during this writing. A partial review
is described by its actual scope. A source transcription, a
successful typesetting run, and a literature search do not create
a mathematical proof check.
